\loai{Loại 1: Chuyển động Brown}
\begin{ex}%chuan
	Trong thí nghiệm của Brown, các hạt phấn hoa chuyển động hỗn độn không ngừng vì
	
	\choice
	{giữa chúng có khoảng cách}
	{chúng là các phân tử}
	{\True các phân tử nước chuyển động không ngừng, va chạm vào chúng từ mọi phía}
	{chúng là các thực thể sống}
	\loigiai{
		\begin{itemize}
			\item Trong thí nghiệm của Brown, sở dĩ các hạt phấn hoa lại chuyển động hỗn độn không ngừng là vì các phân tử nước chuyển động không ngừng va chạm vào các hạt phấn hoa từ mọi phía.
		\end{itemize}
	}
	\haidong{3}\end{ex}
\loai{Loại 2: Cấu tạo chất}
\begin{ex}%chuan
	Nguyên tử, phân tử \textbf{không} có tính chất nào sau đây?
	\choice
	{Chuyển động không ngừng}
	{Giữa chúng có khoảng cách}
	{\True Nở ra khi nhiệt độ tăng, co lại khi nhiệt độ giảm}
	{Chuyển động càng nhanh khi nhiệt độ càng cao}
	\loigiai{
		\begin{itemize}
			\item Nguyên tử, phân tử không có tính chất nở ra khi nhiệt độ tăng, co lại khi nhiệt độ giảm, chỉ có chuyển động càng nhanh khi nhiệt độ càng cao
		\end{itemize}
	}
	\haidong{2}\end{ex}
\loai{Loại 3: Mô hình động học phân tử}
\begin{ex}%chuan
	Theo mô hình động học phân tử thì các chất
	\choice
	{\True được cấu tạo từ các hạt riêng biệt vô cùng nhỏ bé gọi là các phân tử}
	{liền một khối}
	{không được cấu tạo từ các hạt riêng biệt}
	{liền một khối, không được cấu tạo từ các hạt riêng biệt}
	\loigiai{
		\begin{itemize}
			\item Các chất được cấu tạo từ các hạt riêng biệt là các phân tử.
		\end{itemize}
		
	}
	\haidong{2}\end{ex}
	\begin{ex}%chuan \nguon{4}
		Theo mô hình động học phân tử, nhiệt độ của vật cao hay thấp là do
		\choice
		{\True sự chuyển động nhanh hay chậm của các phân tử cấu tạo nên vật}
		{số phân tử cấu tạo nên vật là nhiều hay ít}
		{mật độ các phân tử trên một đơn vị thể tích là lớn hay nhỏ}
		{khối lượng của các phân tử là nặng hay nhẹ}
		\loigiai{
			\begin{itemize}
				\item Nhiệt độ của vật được xác định bởi sự chuyển động nhanh hay chậm của các phân tử cấu tạo nên vật.
				\item Chuyển động của các phân tử nhanh hơn tương ứng với nhiệt độ cao hơn, trong khi chuyển động chậm hơn tương ứng với nhiệt độ thấp hơn.
			\end{itemize}
		}
	\end{ex}
	\begin{ex}%chuan
		Tính chất nào sau đây \textbf{không} phải là của phân tử?
		\choice
		{Chuyển động không ngừng}
		{Giữa các phân tử có khoảng cách}
		{\True Có lúc đứng yên, có lúc chuyển động}
		{Chuyển động càng nhanh thì nhiệt độ của vật càng cao}
		\loigiai{
			\begin{itemize}
				\item Các nguyên tử, phân tử chuyển động không ngừng
			\end{itemize}
		}
		\haidong{2}\end{ex}
\loai{Loại 4: Lực tương tác phân tử}
\begin{ex}%chuan
	Phát biểu nào sau đây là đúng khi nói về lực tương tác giữa các phân tử?
	\choice
	{\True Giữa các phân tử đồng thời có cả lực hút và lực đẩy}
	{Giữa các phân tử chỉ có lực hút hoặc lực đẩy}
	{Giữa các phân tử chỉ có lực đẩy}
	{Giữa các phân tử chỉ có lực hút}
	\loigiai{
		\begin{itemize}
			\item Giữa các phân tử có lực hút và lực đẩy, gọi chung là lực liên kết phân tử.
		\end{itemize}
	}
	\haidong{2}\end{ex}
	\begin{ex}%chuan %[1P5Y1-1][Loai1]
		Khi khoảng cách giữa các phân tử rất nhỏ so với kích thước phân tử, thì giữa các phân tử
		\choice
		{có cả lực hút và lực đẩy, nhưng lực đẩy nhỏ hơn lực hút}
		{\True có cả lực hút và lực đẩy, nhưng lực đẩy lớn hơn lực hút}
		{chỉ có lực đẩy}
		{chỉ có lực hút}	
		\loigiai{
			\begin{itemize}
				\item Khi khoảng cách giữa các phân tử rất nhỏ, thì giữa các phân tử có cả lực hút và lực đẩy, nhưng lực đẩy lớn hơn lực hút.
			\end{itemize}
		}%<MyLT1>
		\haidong{2}\end{ex}
\loai{Loại 5. Chuyển động nhiệt của phân tử}
\begin{ex}%chuan \nguon{1}
	Tốc độ chuyển động nhiệt của các phân tử có ảnh hưởng trực tiếp đến đại lượng nào sau đây?
	\choice
	{\True Nhiệt độ của vật}
	{Trọng lượng riêng của vật}
	{Khối lượng của vật}
	{Thể tích của vật}
	\loigiai{
		\begin{itemize}
			\item Vận tốc chuyển động của các phân tử trong một chất có liên quan mật thiết đến nhiệt độ của chất đó. Khi nhiệt độ tăng, năng lượng nhiệt của các phân tử tăng lên, khiến chúng chuyển động nhanh hơn. Ngược lại, khi nhiệt độ giảm, vận tốc chuyển động của các phân tử cũng giảm theo.		
			\item Các đại lượng khác như trọng lượng riêng, khối lượng, và thể tích của vật không ảnh hưởng trực tiếp đến vận tốc chuyển động của các phân tử. Do đó, đáp án đúng là nhiệt độ của vật.
		\end{itemize}
	}
\end{ex}
\loai{Loại 6: Xác định các thể}
\begin{ex}%chuan
	Phát biểu nào sau đây là \textbf{sai}? Cùng một khối lượng của một chất nhưng khi ở các thể khác nhau thì sẽ khác nhau về
	\choice
	{thể tích}
	{khối lượng riêng}
	{\True kích thước của các nguyên tử}
	{trật tự của các nguyên tử}
	
	\loigiai{
		\begin{itemize}
			\item Sử dụng lý thuyết các thể của chất: Kích thước của các nguyên tử ở các thể khác nhau là như nhau
		\end{itemize}
	}
	\haidong{1}\end{ex}
	\begin{ex}%chuan
		Cho các chất sau: đồng (copper); nitrogen; vàng (gold); nước; oxygen; dầu hỏa; sắt (iron). Số chất ở thể rắn trong điều kiện thường là
		
		\choice
		{1}
		{2}
		{\True 3}
		{4}
		\loigiai{
			
			\begin{itemize}
				\item Các chất ở thể rắn trong điều kiện thường là: đồng (copper); vàng (gold); sắt (iron).
			\end{itemize}
			
		}
		\haidong{2}\end{ex}
\loai{Loại 7: Cấu trúc chất rắn}
\begin{ex}%chuan
	Hình dưới đây biểu diễn mô hình cấu tạo phân tử của ba chất $A,\ B$ và $C$.
	\begin{center}
		\begin{tikzpicture}[scale=1.1,draw=Mapcolor, very thick, >=stealth']
			\tkzDefPoints{0.2/2/a, .8/2/b, 3.2/0/c, 3.5/.45/d}
			\draw (-0.5,-0.5) -- (4,-0.5) -- (4,2.5) -- (-0.5,2.5) -- (-0.5,-0.5);
			\node [draw, thick, shape=circle, minimum width=.65cm, minimum height=.5cm, anchor=center] at (b) {};
			\node [draw, thick,fill=white, shape=circle, minimum width=.65cm, minimum height=.5cm, anchor=center] at (a) {};
			\node [draw, thick,shape=circle, minimum width=.65cm, minimum height=.5cm, anchor=center] at (d) {};
			\node [draw, thick,fill=white, shape=circle, minimum width=.65cm, minimum height=.5cm, anchor=center] at (c) {};
			\node at (1.75,-1) {$A$};
			
			\begin{scope}[shift={(5,0)},rotate=0]
				\tkzDefPoints{0.5/2/a, 1.5/2/b, 2.5/2/c, 3.5/2/d}
				\tkzDefPoints{0.5/1/x, 1.5/1/y, 2.5/1/z, 3.5/1/o}
				\tkzDefPoints{0.5/0/k, 1.5/0/m, 2.5/0/n, 3.5/0/p}
				
				\draw (0,-0.5) -- (4,-0.5) -- (4,2.5) -- (0,2.5) -- (0,-0.5);
				\node [draw, thick, shape=circle, minimum width=.95cm, minimum height=.95cm, anchor=center] at (a) {};
				\node [draw, thick, shape=circle, minimum width=.95cm, minimum height=.95cm, anchor=center] at (b) {};
				\node [draw, thick, shape=circle, minimum width=.95cm, minimum height=.95cm, anchor=center] at (c) {};
				\node [draw, thick, shape=circle, minimum width=.95cm, minimum height=.95cm, anchor=center] at (d) {};
				
				\node [draw, thick, shape=circle, minimum width=.95cm, minimum height=.95cm, anchor=center] at (x) {};
				\node [draw, thick, shape=circle, minimum width=.95cm, minimum height=.95cm, anchor=center] at (y) {};
				\node [draw, thick, shape=circle, minimum width=.95cm, minimum height=.95cm, anchor=center] at (z) {};
				\node [draw, thick, shape=circle, minimum width=.95cm, minimum height=.95cm, anchor=center] at (o) {};
				
				\node [draw, thick, shape=circle, minimum width=.95cm, minimum height=.95cm, anchor=center] at (k) {};
				\node [draw, thick, shape=circle, minimum width=.95cm, minimum height=.95cm, anchor=center] at (m) {};
				\node [draw, thick, shape=circle, minimum width=.95cm, minimum height=.95cm, anchor=center] at (n) {};
				\node [draw, thick, shape=circle, minimum width=.95cm, minimum height=.95cm, anchor=center] at (p) {};
				
				\node at (2,-1) {$B$};
			\end{scope}
			
			\begin{scope}[shift={(10,0)},rotate=0]
				\tkzDefPoints{0.5/2/a, 1.25/1.65/b, 2.5/1.3/c, 2.3/1.75/d}
				\tkzDefPoints{0.75/1.5/x, 3.5/2/y, 4.15/1.5/z, 3.65/1.5/o}
				\tkzDefPoints{1.5/.5/k, 1/0.25/m, 2.75/0/n, 3.5/0.5/p, 3/0.45/q}
				
				\draw (0,-0.5) -- (4.5,-0.5) -- (4.5,2.5) -- (0,2.5) -- (0,-0.5);
				\node [draw, thick, shape=circle, minimum width=.6cm, minimum height=.5cm, anchor=center] at (b) {};
				\node [draw, thick,fill=white, shape=circle, minimum width=.6cm, minimum height=.5cm, anchor=center] at (a) {};
				\node [draw, thick,fill=white, shape=circle, minimum width=.8cm, minimum height=.5cm, anchor=center] at (x) {};
				
				\node [draw, thick, shape=circle, minimum width=.6cm, minimum height=.5cm, anchor=center] at (c) {};
				\node [draw, thick,fill=blue,shape=circle, minimum width=.9cm, minimum height=.5cm, anchor=center] at (d) {};
				
				\node [draw, thick, shape=circle, minimum width=.6cm, minimum height=.5cm, anchor=center] at (y) {};
				\node [draw, thick,fill=white, shape=circle, minimum width=.6cm, minimum height=.5cm, anchor=center] at (z) {};
				\node [draw, thick,fill=white, shape=circle, minimum width=.8cm, minimum height=.5cm, anchor=center] at (o) {};
				
				\node [draw, thick, shape=circle, minimum width=.6cm, minimum height=.5cm, anchor=center] at (k) {};
				\node [draw, thick,fill=blue,shape=circle, minimum width=.9cm, minimum height=.5cm, anchor=center] at (m) {};
				
				\node [draw, thick, shape=circle, minimum width=.6cm, minimum height=.5cm, anchor=center] at (n) {};
				\node [draw, thick,fill=white, shape=circle, minimum width=.6cm, minimum height=.5cm, anchor=center] at (p) {};
				\node [draw, thick,fill=white, shape=circle, minimum width=.8cm, minimum height=.5cm, anchor=center] at (q) {};
				\node at (2.25,-1) {$C$};
			\end{scope}
			
			\node at (7,-2) {};
		\end{tikzpicture}
	\end{center}
	\vspace*{-1.5 cm}
	Mô hình chất rắn tinh thể giống với mô hình nào nhất sau đây?
	\choice
	{A}
	{\True B}
	{C}
	{A và B}
	\loigiai{
		\begin{itemize}
			\item Từ ba mô hình đã cho, chất B là chất rắn vì các phân tử ở gần nhau nhất và được sắp xếp theo trật tự xác định.
		\end{itemize}
		
	}
	
\end{ex}
\begin{ex}%chuan
	Đặc điểm nào sau đây là của chất rắn?
	\choice
	{\True Lực tương tác giữa các phân tử rất mạnh}
	{Không có hình dạng xác định}
	{Lực tương tác giữa các phân tử rất yếu}
	{Không có thể tích riêng xác định}
	\loigiai{
		\begin{itemize}
			\item Trong chất rắn, các phân tử ở rất gần nhau nên lực tương tác giữa các phân tử chất rắn rất mạnh.
		\end{itemize}
	}
	\haidong{2}\end{ex}
	\begin{ex}%chuan
		Phát biểu nào dưới đây là \textbf{không} đúng với thể rắn?
		\choice
		{Các phân tử rất gần nhau (cỡ kích thước phân tử)}
		{Các phân tử sắp xếp có trật tự}
		{Các phân tử dao động quanh vị trí cân bằng cố định}
		{\True Các phân tử dao động quanh vị trí cân bằng không xác định}
		\loigiai{
			\begin{itemize}
				\item Các phân tử chất rắn được giữ ở các vị trí cân bằng và mỗi phân tử
				chỉ có thể dao động xung quanh vị trí cân bằng xác định này.
			\end{itemize}
		}
		\haidong{2}\end{ex}
\loai{Loại 8: Chất rắn kết tinh}
\begin{ex}%chuan
	Các nguyên tử sắp xếp theo một trật tự nhất định trong không gian được gọi là
	\choice
	{mạng nguyên tử}
	{mạng internet}
	{\True mạng tinh thể}
	{mạng nhện}
	\loigiai{
		\begin{itemize}
			\item Các phân tử liên kết chặt với nhau và sắp xếp theo một trật tự hình học xác
			định, tuần hoàn trong không gian, gọi là mạng tinh thể.
		\end{itemize}
	}
	\haidong{2}\end{ex}
	\begin{ex}%chuan%[1P7B1-1][Loai1]
		Chất nào sau đây \textbf{không} có cấu trúc tinh thể?
		\choice
		{\True Miếng nhựa thông}
		{Hạt đường}
		{Viên kim cương}
		{Khối thạch anh}
		\loigiai{
			\begin{itemize}
				\item Miếng nhựa thông là chất rắn vô định hình, không có cấu trúc tinh thể
			\end{itemize}
		}	
		%<MyLT3>
		\haidong{2}\end{ex}
		\begin{ex}%chuan \nguon{1}
			\immini[thm]{Hình bên là cấu trúc tinh thể của
				\motcot
				{\True muối ăn}
				{kim cương}
				{thủy tinh}
				{đồng}}
			{\includegraphics[width=3cm]{images/2024_05_26_010fc85c05a422f9b1ffg-05}} 
			\loigiai{
				Muối ăn có cấu trúc lập phương.
			}
		\end{ex}
		
\loai{Loại 9: Chất rắn vô định hình}
\begin{ex}%chuan%[1P7B1-1][Loai2]
	Chất rắn nào dưới đây là chất rắn vô định hình?
	\choice
	{\True Thủy tinh}
	{Băng phiến}
	{Hợp kim}
	{Kim loại}
	\loigiai{
		\begin{itemize}
			\item Các chất rắn vô định hình như thủy tinh, các loại nhựa, cao su, ...
		\end{itemize}
	}%<MyLT1>
	\haidong{3}\end{ex}
\loai{Loại 10: Chất rắn kết tinh – Chất rắn vô định hình}

	\begin{ex}%chuan
		Chất rắn vô định hình là chất rắn
		\choice
		{có cấu trúc tinh thể}
		{\True không có cấu trúc tinh thể}
		{có cấu trúc mạng lập phương}
		{có cấu trúc mạng lục giác}
		\loigiai{
			\begin{itemize}
				\item Chất rắn vô định hình: Không có cấu trúc tinh thể.
				\item Thuỷ tinh, nhựa đường, cao su,... là những chất rắn vô định hình.
			\end{itemize}
		}
		\haidong{2}\end{ex}
	\begin{ex}%chuan
		Chất rắn nào dưới đây \textbf{không} phải là chất rắn vô định hình?
		\choice
		{\True Thạch anh}
		{Thủy tinh}
		{Sáp}
		{Cao su}
		\loigiai{
			\begin{itemize}
				\item Thạch anh có cấu trúc tinh thể, là chất rắn kết tinh
			\end{itemize}
		}
		\haidong{3}\end{ex}
\loai{Loại 11: Cấu trúc chất lỏng}
\begin{ex}%chuan
	Đặc điểm nào sau đây là đặc điểm của thể lỏng?	
	\choice
	{Khoảng cách giữa các phân tử rất lớn so với kích thước của chúng}
	{\True Lực tương tác phân tử yếu hơn lực tương tác phân tử ở thể rắn}
	{Không có thể tích và hình dạng riêng xác định}
	{Các phân tử dao động xung quanh vị trí cân bằng xác định}
	\loigiai{
		\begin{itemize}
			\item Khoảng chất lỏng, các phân tử ở xa nhau hơn so với các phân	tử trong chất rắn nên lực tương tác giữa các phân tử chất lỏng nhỏ hơn trong chất
			rắn
		\end{itemize}
		
	}
	\haidong{3}\end{ex}
	\begin{ex}%chuan
		Trong chuyển động nhiệt, các phân tử của chất lỏng
		\choice
		{chuyển động hoàn toàn hỗn loạn}
		{chuyển động hỗn loạn quanh vị trí cân bằng xác định}
		{dao động xung quanh các vị trí cân bằng cố định}
		{\True dao động quanh vị trí cân bằng nhưng những vị trí cân bằng này không cố định mà di chuyển}
		\loigiai{
			\begin{itemize}
				\item Các phân tử của chất lỏng dao động xung quanh các vị trí cân bằng và các vị trí cân bằng này lại có thể
				thay đổi.
			\end{itemize}
		}
		\haidong{2}\end{ex}
		\begin{ex}%chuan \nguon{4}
			Chất lỏng \textbf{không} có hình dạng xác định vì các phân tử chất lỏng
			\choice
			{dao động tại các vị trí cân bằng xác định}
			{có thể chuyển động phân tán ra xa nhau}
			{\True dao động quanh các vị trí cân bằng có thể dịch chuyển được}
			{có thể chuyển động tự do}
			\loigiai{
				Chất lỏng có tính linh động vì các phân tử của nó dao động quanh các vị trí cân bằng có thể di chuyển được, giúp chất lỏng dễ dàng thay đổi hình dạng theo hình dạng của vật chứa.
			}
		\end{ex}
		\begin{ex}%chuan \nguon{1}
			Khi đổ 50 lít rượu vào 50 lít nước thì thu được bao nhiêu lít hỗn hợp?
			\choice
			{100 lít}
			{50 lít}
			{Lớn hơn 100 lít}
			{\True Nhỏ hơn 100 lít}
			\loigiai{
				\begin{itemize}
					\item Khi đổ 50 lít rượu vào 50 lít nước thì thu được \textbf{nhỏ hơn 100 lít} hỗn hợp. 
					\item Nguyên nhân là do các phân tử rượu và nước có thể xen kẽ và chen vào các khoảng trống giữa chúng, làm thể tích tổng giảm so với khi đo riêng rẽ từng phần. 
					\item Đây là một hiện tượng phổ biến khi trộn các chất có liên kết phân tử khác nhau.\end{itemize}
			}
		\end{ex}
		\begin{ex}%chuan
			Hình bên là hình ảnh phóng to của bề mặt bàn. Chất lỏng di chuyển trên mặt bàn dễ dàng hơn so với chất rắn do
				\immini{\choice
				{các phân tử trong chất lỏng có thể chuyển động tự do tương đối, giúp chúng chảy và thay đổi hình dạng}
				{chất lỏng có khối lượng nhỏ hơn chất rắn nên ít chịu tác động của trọng lực}
				{lực hút giữa các phân tử trong chất lỏng lớn hơn so với chất rắn nên chất lỏng trượt dễ dàng hơn}
				{chất lỏng có thể bị nén nhiều hơn so với chất rắn, làm giảm ma sát với mặt bàn}}{	\includegraphics[width=6 cm]{images/caibannew}}
			\loigiai{
				\begin{itemize}
					\item Chất lỏng có thể di chuyển trên mặt bàn dễ dàng hơn so với chất rắn vì trong chất lỏng, các phân tử có sự tự do chuyển động tương đối, cho phép chúng chảy và thay đổi hình dạng. Ngược lại, các phân tử trong chất rắn được sắp xếp cứng nhắc và cố định tại các vị trí nhất định, do đó chất rắn không thể di chuyển hay chảy dễ dàng như chất lỏng.
				\end{itemize}
			}
		\end{ex}
\loai{Loại 12: Cấu trúc chất khí}
\begin{ex}%chuan %[1P5Y1-1][Loai1]
	Nguyên nhân gây ra chuyển động hoàn toàn hỗn loạn của các phân tử chất khí là do
	\choice
	{các phân tử luôn va chạm với nhau}
	{\True lực tương tác giữa các phân tử rất yếu}
	{khoảng cách giữa các phân tử nhỏ}
	{khoảng cách giữa các phân tử không giống nhau}
	\loigiai{
		\begin{itemize}
			\item Ở thể khí, lực tương tác giữa các phân tử rất yếu nên các phân tử chuyển động hoàn toàn hỗn loạn.
		\end{itemize}
	}%<MyLT2>
	\haidong{2}\end{ex}
	\begin{ex}%chuan%[1P5Y1-1][Loai3]
		Phát biểu nào sau đây là \textbf{sai} khi nói về chất khí?
		\choice
		{Lực tương tác giữa các phân tử là rất yếu}
		{\True Các phân tử khí ở rất gần nhau}
		{Chất khí không có hình dạng và thể tích riêng}
		{Chất khí luôn chiếm toàn bộ thể tích bình chứa và nén được dễ dàng}
		\loigiai{
			\begin{itemize}
				\item Phát biểu “Các phân tử khí ở rất gần nhau” là sai. Thực ra, các phân tử khí ở rất xa nhau (khoảng cách giữa các phân tử lớn gấp hàng chục lần kích thước của chúng).
			\end{itemize}
		}%<MyLT1>
		\haidong{3}\end{ex}
		\begin{ex}%chuan
			Chất khí luôn luôn chiếm hết thể tích của bình chứa bởi vì
			
			\choice
			{\True lực liên kết giữa các phân tử chất khí rất yếu nên chúng chuyển động tự do về mọi phía}
			{các phân tử chất khí luôn chuyển động nhiệt không ngừng}
			{trong chất khí có quá nhiều phân tử}
			{các phân tử chất khí luôn luôn đẩy nhau ra xa nên chúng cách nhau càng xa càng tốt}
			\loigiai{
				\begin{itemize}
					\item  Khoảng cách giữa các phân tử rất lớn so với kích
					thước của chúng nên lực tương tác giữa các phân tử hầu như không đáng
					kể (trừ khi va chạm nhau) nên chúng chuyển động tự do về mọi phía
				\end{itemize}
			}
			\haidong{3}\end{ex}
			\begin{ex}%chuan
				\immini[thm]{Hình bên mô tả cấu trúc phân tử ở thể nào dưới đây?
					\choice[2]
					{Thể lỏng}
					{\True Thể khí}
					{Thể rắn}
					{Plasma}}
				{\begin{tikzpicture}[scale=.8,thick,>=stealth']
						\draw[fill=blue!10] (0,0) -- (0,-2) -- (5,-2) -- (5,0) -- (0,0);
						\shade[ball color=blue!70] (0.4,-0.4) circle (1.3ex);
						%\shade[ball color=blue!70] (1,-1.5) circle (1.3ex);
						\shade[ball color=blue!70] (2,-1.6) circle (1.3ex);
						\shade[ball color=blue!70] (3.2,-0.3) circle (1.3ex);
						%\shade[ball color=blue!70] (4.25,-0.75) circle (1.3ex);
						\shade[ball color=blue!70] (4.5,-1.5) circle (1.3ex);
						\shade[ball color=blue!70] (5.5,-0.25) circle (1.3ex);
						\shade[ball color=blue!30] (6,-1) circle (1.3ex);
						\shade[ball color=blue!30] (5.5,-1.75) circle (1.3ex);
						\draw[dashed,thin,->] (5.6,-0.4) -- (5.9,-0.8);
						\draw [dashed,thin,->] (6,-1) -- (5.6,-1.6);
						\draw [dashed,thin,->] (5.5,-1.8) -- (6.1,-2);
				\end{tikzpicture}}
				\loigiai{
					\begin{itemize}
						\item  Khoảng cách giữa các phân tử rất lớn nên đây là thể khí
					\end{itemize}
				}
				\haidong{2}\end{ex}
\loai{Loại 13: Tổng hợp các thể vật chất đơn giản}
\begin{ex}%chuan%[1P5Y1-1][Loai1]
	Chất nào rất khó nén được?
	\choice
	{Chỉ có chất rắn}
	{\True Chất rắn, chất lỏng}
	{Chất khí, chất lỏng}
	{Chất khí, chất rắn}
	\loigiai{
		\begin{itemize}
			\item Do khoảng cách các phân tử ở gần nhau nên chất lỏng và chất rắn gần như không nén được.
		\end{itemize}
	} 
	\haidong{2}\end{ex}
\loai{Loại 14: So sánh đặc điểm của các chất}
\begin{ex}%chuan \nguon{1}
	Kim cương và than chì có tính chất vật lí khác nhau là do chúng khác nhau về
	\choice
	{\True cấu trúc mạng tinh thể}
	{bản chất các hạt tạo thành tinh thể}
	{loại liên kết giữa các hạt trong tinh thể}
	{kích thước tinh thể}
	\loigiai{
		\begin{itemize}
			\item Kim cương và than chì đều được tạo thành từ các nguyên tử cacbon, nhưng chúng có cấu trúc tinh thể khác nhau. 
			\item Trong kim cương, mỗi nguyên tử carbon liên kết với bốn nguyên tử carbon khác xung quanh theo cấu trúc tứ diện, tạo nên một mạng lưới ba chiều rất cứng và bền. 
			\item Trong khi đó, trong than chì, các nguyên tử carbon liên kết tạo thành tấm phẳng với các liên kết mạnh trong từng lớp, nhưng các lớp này lại liên kết với nhau bằng các lực yếu hơn (lực Van der Waals), cho phép chúng trượt lên nhau dễ dàng. 
			\item Sự khác biệt trong cấu trúc tinh thể này dẫn đến các tính chất vật lí rất khác nhau giữa kim cương và than chì.\end{itemize}
	}
\end{ex}
\begin{ex}%chuan \nguon{4}
	Khi nói về khoảng cách trung bình giữa các phân tử trong chất rắn, chất lỏng, chất khí. Phát biểu nào sau đây là đúng?
	\choice
	{Khoảng cách giữa các phân tử trong chất lỏng xa hơn so với các phân tử trong chất khí}
	{Khoảng cách giữa các phân tử trong chất rắn xa hơn so với các phân tử trong chất lỏng}
	{\True Khoảng cách giữa các phân tử trong chất lỏng gần hơn so với các phân tử trong chất khí}
	{Khoảng cách giữa các phân tử trong chất khí gần hơn so với các phân tử trong chất lỏng}
	\loigiai{
		\begin{itemize}
			\item Trong chất lỏng, các phân tử ở gần nhau hơn so với trong chất khí nhưng xa nhau hơn so với trong chất rắn.
			\item Chất khí có khoảng cách giữa các phân tử lớn nhất, tạo nên tính chất dễ nén.
			\item Chất rắn có khoảng cách giữa các phân tử nhỏ nhất, tạo nên cấu trúc vững chắc.
		\end{itemize}
	}
\end{ex}
\begin{ex}%chuan \nguon{1}
	Chất lỏng được dẫn bằng các đường ống. Điều này là do yếu tố nào sau đây?
	\choice
	{\True Tính chất thay đổi hình dạng theo bình chứa}
	{Chất lỏng có cấu trúc mạng tinh thể}
	{Các phân tử chất lỏng chuyển động hỗn loạn} 
	{Các phân tử chất lỏng va chạm với thành bình}
	\loigiai{
		Thể lỏng có hình dạng thay đổi được theo hình dạng của phần bình chứa nó.
	}
\end{ex}
\loai{Loại 16: Câu hỏi đúng sai}
\begin{ex}%chuan \nguon{4}
	Theo mô hình động học phân tử thì các phân tử cấu tạo nên các vật có sự tương tác với nhau. Sự khác biệt về độ lớn của lực tương tác giữa các phân tử trong chất rắn, chất lỏng, chất khí dẫn đến sự khác nhau về cấu trúc của chúng.
	\choiceTF[t]
	{\True Lực tương tác giữa các phân tử chất rắn mạnh hơn lực tương tác giữa các phân tử chất khí}
	{\True Lực tương tác giữa các phân tử chất rắn rất mạnh nên giữ các phân tử chỉ dao động quanh các vị trí cân bằng cố định, vì vậy chất rắn có hình dạng xác định}
	{Các phân tử chất lỏng tương tác với nhau rất yếu nên các phân tử chất lỏng có thể chuyển động tự do ra xa nhau, vì vậy chất lỏng không có hình dạng xác định}
	{Các chất khí có khoảng cách giữa các phân tử rất lớn so với kích thước của chúng nên lực tương tác giữa các phân tử gần như không đáng kể}
	\loigiai{
		\begin{itemchoice}
			\itemch 
			\itemch 
			\itemch Lực tương tác giữa các phân tử chất lỏng vẫn đủ để giữ các phân tử chất lỏng không chuyển động phân tán ra xa nhau. 
			\itemch 
		\end{itemchoice}
	}
\end{ex}
\begin{ex}%chuan \nguon{1}
	\immini[thm]{Chuyển động hỗn loạn, không ngừng của các hạt rất nhỏ dưới tác dụng của các phân tử chất lỏng hoặc chất khí được gọi là chuyển động Brown.}
	{\includegraphics[width=3cm]{images/2024_05_26_010fc85c05a422f9b1ffg-01(1)}}
	\choiceTF[t]
	{Chuyển động Brown do một nhà Vật Lí học phát hiện ra}
	{Khi quan sát các hạt phấn hoa, Brown khẳng định cấu tạo chất bao gồm các phân tử}
	{\True Nguyên nhân của chuyển động Brown là do các phân tử nước luôn luôn chuyển động hỗn loạn, va chạm vào các hạt phấn hoa từ nhiều phía}
	{Hạt phấn hoa chuyển động có quy luật xác định}	
	\loigiai{
		\begin{itemchoice}
			\itemch 
			\itemch 
			\itemch 
			\itemch 
		\end{itemchoice} 
	}
\end{ex}

\begin{ex}%chuan \nguon{4}
	So sánh giữa nước ở dạng lỏng và nước ở dạng hơi (hơi nước) với cùng khối lượng theo mô hình động học phân tử.
	\choiceTF[t]
	{Hơi nước có khối lượng riêng lớn hơn nước lỏng}
	{\True Hơi nước có thể tích lớn hơn nước lỏng}
	{\True Trong hơi nước, các phân tử nước có khoảng cách xa nhau hơn so với khi ở dạng nước lỏng}
	{Hơi nước và nước lỏng đều có các phân tử di chuyển hỗn loạn}
	\loigiai{
		\begin{itemchoice}
			\itemch 
			\itemch 
			\itemch 
			\itemch 
		\end{itemchoice}
		Theo mô hình động học phân tử, các chất ở thể khí có khoảng cách giữa các phân tử xa hơn khoảng cách giữa các phân tử khi nó ở thể lỏng. 
		Vì vậy với cùng một khối lượng chất như nhau, thể tích khi ở thể khí lớn hơn thể lỏng. 
		Điều này làm cho khối lượng riêng của một chất ở thể khí nhỏ hơn khối lượng riêng của chất đó khi ở thể lỏng.
	}
\end{ex}

